\section{Universality}

The mesh is not only a stage for physics, it is a computer, and the same growing weave that does physics computes. One
local rule is enough to compute anything at all.

A computer needs only a few things: signals that travel, signals that interact in a way rich enough to build a logic
gate, and a way to hold state. The mesh has all three. A free tone's value is handed one dock per beat along its
direction, a simple signal that propagates. Tones that meet collide by the knit, and the collision is rich enough that the right
arrangement of crossing tones realizes the basic logical operations, gliders that meet and interact like the parts of a
circuit. And because the rule is reversible, it loses no information, which is exactly the basis for reversible
computing, where a computation can in principle be run backward and wastes no information.

So the reversible local rule supports glider-like signals and interactions rich enough to be computationally universal,
able to simulate any computer, which is the strongest statement one can make about a substrate's computational power.
This is not an extra mechanism bolted on, it is the same knit that moves the tones, read as logic instead of as physics.
The universe and a universal computer are the same object seen two ways, the physical reading and the computational
reading of one growing rule.

This matters for the larger story, because a mind, later, is a region of mesh that computes, and the room and the
speed it needs come from the same hyperbolic geometry that gave holography and the mass hierarchy. The next sections show
the two practical things that make the mesh not just universal but usable: a fast address for every cell, and a
natural place to store and recall patterns.
